\input{preamble}

% OK, start here.
%
\begin{document}

\title{Retour sur la résolution des surfaces}


\maketitle

\phantomsection
\label{section-phantom}

\tableofcontents

\section{Introduction}
\label{section-introduction}

\noindent
Ce chapitre traite de la résolution des singularités des
espaces algébriques noethériens de dimension $2$.
Dans un chapitre antérieur, nous avons déjà étudié, à la suite de
Lipman \cite{Lipman}, la résolution des surfaces dans le cas des schémas.
Voir Résolution des surfaces, section
\ref{resolve-section-introduction}.
La plupart des résultats de ce chapitre se déduisent immédiatement
des résultats relatifs aux schémas.

\medskip\noindent
Sauf mention expresse du contraire, tous les objets géométriques
de ce chapitre seront des espaces algébriques. Ainsi, lorsque nous dirons
``soit $f : X \to Y$ une modification'', cela signifiera que
$f$ est un morphisme au sens de Espaces au-dessus de corps, définition
\ref{spaces-over-fields-definition-modification}.
Il en ira de même pour les morphismes propres, et ainsi de suite.










\section{Modifications}
\label{section-modifications}

\noindent
Soit $(A, \mathfrak m, \kappa)$ un anneau local noethérien. Posons
$S = \Spec(A)$ et $U = S \setminus \{\mathfrak m\}$. Dans cette section,
nous considérerons la catégorie
\begin{equation}
\label{equation-modification}
\left\{
f : X \longrightarrow S
\quad \middle| \quad
\begin{matrix}
X\text{ est un espace algébrique}\\
f\text{ est un morphisme propre}\\
f^{-1}(U) \to U\text{ est un isomorphisme}
\end{matrix}
\right\}
\end{equation}
Un morphisme de $X/S$ vers $X'/S$ sera un morphisme d'espaces algébriques
$X \to X'$ compatible aux morphismes structuraux au-dessus de $S$. Dans
Algébrisation des espaces formels, section
\ref{restricted-section-modifications}, nous avons vu que cette catégorie
ne dépend que du complété de $A$ et démontré quelques propriétés élémentaires
de ses objets. Dans cette section, nous étudions plus particulièrement les cas
où $\dim(A) \leq 2$ ou bien où la fibre fermée est de dimension au plus $1$.

\begin{lemma}
\label{lemma-modification}
Soit $(A, \mathfrak m, \kappa)$ un domaine local noethérien
de dimension $2$ tel que $U = \Spec(A) \setminus \{\mathfrak m\}$
soit un schéma normal. Alors toute modification $f : X \to \Spec(A)$
est un morphisme comme dans (\ref{equation-modification}).
\end{lemma}

\begin{proof}
Soit $f : X \to S$ une modification. Il nous faut montrer que
$f^{-1}(U) \to U$ est un isomorphisme. Comme tout point fermé $u$ de $U$
est de codimension $1$, cela résulte de
Espaces au-dessus de corps, lemme
\ref{spaces-over-fields-lemma-modification-normal-iso-over-codimension-1}.
\end{proof}

\begin{lemma}
\label{lemma-closed-immersion-on-fibre}
Soit $(A, \mathfrak m, \kappa)$ un anneau local noethérien.
Soit $g : X \to Y$ un morphisme de la catégorie
(\ref{equation-modification}). Si le morphisme induit
$X_\kappa \to Y_\kappa$ entre les fibres spéciales est une immersion fermée,
alors $g$ est une immersion fermée.
\end{lemma}

\begin{proof}
C'est un cas particulier de
Compléments sur les morphismes d'espaces, lemme
\ref{spaces-more-morphisms-lemma-where-closed-immersion}.
\end{proof}

\begin{lemma}
\label{lemma-dimension-special-fibre}
Soit $(A, \mathfrak m, \kappa)$ un domaine local noethérien
de dimension $\geq 1$.
Soit $f : X \to \Spec(A)$ un morphisme d'espaces algébriques.
Supposons satisfaite au moins l'une des conditions suivantes :
\begin{enumerate}
\item $f$ est une modification (Espaces au-dessus de corps, définition
\ref{spaces-over-fields-definition-modification}),
\item $f$ est une altération (Espaces au-dessus de corps, définition
\ref{spaces-over-fields-definition-alteration}),
\item $f$ est localement de type fini et quasi-séparé, $X$ est intègre,
et il existe exactement un point de $|X|$ dont l'image est le point générique
de $\Spec(A)$,
\item $f$ est localement de type fini, $X$ est raisonnable, et les points
de $|X|$ dont l'image est le point générique de $\Spec(A)$ sont
les points génériques des composantes irréductibles de $|X|$,
\item ajouter d'autres cas ici.
\end{enumerate}
Alors $\dim(X_\kappa) \leq \dim(A) - 1$.
\end{lemma}

\begin{proof}
Les cas (1), (2) et (3) sont des cas particuliers de (4). Choisissons un
schéma affine $U = \Spec(B)$ et un morphisme étale $U \to X$.
L'homomorphisme d'anneaux $A \to B$ est de type fini. Il nous faut montrer que
$\dim(U_\kappa) \leq \dim(A) - 1$. Puisque $X$ est raisonnable, les points
génériques des composantes irréductibles de $U$ sont les points situés
au-dessus des points génériques des composantes irréductibles de $|X|$ ; voir
Espaces raisonnables, lemme
\ref{decent-spaces-lemma-decent-generic-points}.
La fibre de $\Spec(B) \to \Spec(A)$ au-dessus de $(0)$ est donc l'ensemble
(fini) des idéaux premiers minimaux $\mathfrak q_1, \ldots, \mathfrak q_r$
de $B$. Ainsi, $A_f \to B_f$ est fini pour un certain $f \in A$ non nul
(Algèbre, lemme \ref{algebra-lemma-generically-finite}).
Il s'ensuit que $\kappa(\mathfrak q_i)$ est une extension finie du
corps des fractions de $A$.
Soit $\mathfrak q \subset B$ un idéal premier au-dessus de $\mathfrak m$. Alors
$$
\dim(B_\mathfrak q) = \max \dim((B/\mathfrak q_i)_{\mathfrak q})
\leq \dim(A)
$$
l'inégalité résultant de la formule de dimension appliquée à
$A \subset B/\mathfrak q_i$ ; voir Algèbre, lemme
\ref{algebra-lemma-dimension-formula}.
Toutefois, la dimension de $B_\mathfrak q/\mathfrak m B_\mathfrak q$
(qui est l'anneau local de $U_\kappa$ au point correspondant)
est inférieure d'au moins une unité, car les idéaux premiers minimaux
$\mathfrak q_i$ n'appartiennent pas à $V(\mathfrak m)$. On conclut par
Propriétés, lemme \ref{properties-lemma-dimension}.
\end{proof}

\begin{lemma}
\label{lemma-modification-of-dim-2-is-projective-over-complete}
Si $(A, \mathfrak m, \kappa)$ est un domaine local noethérien complet
de dimension $2$, alors toute modification de $\Spec(A)$ est projective sur $A$.
\end{lemma}

\begin{proof}
Par Compléments sur les morphismes d'espaces, lemme
\ref{spaces-more-morphisms-lemma-projective-over-complete}
il suffit de montrer que la fibre spéciale de toute modification
$X$ de $\Spec(A)$ est de dimension $\leq 1$.
Cela résulte du lemme \ref{lemma-dimension-special-fibre}.
\end{proof}






\section{Stratégie}
\label{section-strategy}

\noindent
Soit $S$ un schéma. Soit $X$ un espace algébrique raisonnable sur $S$.
Soient $x_1, \ldots, x_n \in |X|$ des points fermés deux à deux distincts.
Pour chaque $i$, choisissons un voisinage étale élémentaire
$(U_i, u_i) \to (X, x_i)$ comme dans Espaces raisonnables, lemme
\ref{decent-spaces-lemma-decent-space-elementary-etale-neighbourhood}.
Cela signifie que $U_i$ est un schéma affine, que $U_i \to X$ est étale,
que $u_i$ est l'unique point de $U_i$ au-dessus de $x_i$, et que
$\Spec(\kappa(u_i)) \to X$ est un monomorphisme représentant $x_i$.
Quitte à restreindre $U_i$, nous pouvons et voulons supposer que, si
$j \not = i$, aucun point de $U_i$ ne s'envoie sur $x_j$.
Remarquons que $u_i \in U_i$ est un point fermé.

\medskip\noindent
Notons $\mathcal{C}_{X, \{x_1, \ldots, x_n\}}$ la catégorie des
morphismes d'espaces algébriques $f : Y \to X$ qui induisent un isomorphisme
$f^{-1}(X \setminus \{x_1, \ldots, x_n\}) \to
X \setminus \{x_1, \ldots, x_n\}$.
Pour chaque $i$, notons $\mathcal{C}_{U_i, u_i}$ la catégorie des
morphismes d'espaces algébriques $g_i : Y_i \to U_i$ qui induisent un
isomorphisme $g_i^{-1}(U_i \setminus \{u_i\}) \to U_i \setminus \{u_i\}$.
Le changement de base définit un foncteur
\begin{equation}
\label{equation-equivalence}
F :
\mathcal{C}_{X, \{x_1, \ldots, x_n\}}
\longrightarrow
\mathcal{C}_{U_1, u_1} \times \ldots \times \mathcal{C}_{U_n, u_n}
\end{equation}
Le lemme suivant permettra de ramener au cas des schémas au moins certains
des problèmes de ce chapitre.

\begin{lemma}
\label{lemma-equivalence}
Le foncteur $F$ de (\ref{equation-equivalence}) est une équivalence.
\end{lemma}

\begin{proof}
Pour $n = 1$, c'est Limites d'espaces, lemme
\ref{spaces-limits-lemma-excision-modifications}.
Pour $n > 1$, le lemme se démontre exactement de la même manière, ou s'en
déduit. Supposons par exemple que les $g_i : Y_i \to U_i$ soient des objets
de $\mathcal{C}_{U_i, u_i}$. Le cas $n = 1$ fournit alors des
$f'_i : Y'_i \to X$ qui sont des isomorphismes au-dessus de
$X \setminus \{x_i\}$ et dont le changement de base à $U_i$ est $f_i$.
On peut alors poser
$$
f : Y = Y'_1 \times_X \ldots \times_X Y'_n \to X
$$
C'est un objet de $\mathcal{C}_{X, \{x_1, \ldots, x_n\}}$ dont le
changement de base par $U_i \to X$ redonne $g_i$. Le foncteur est donc
essentiellement surjectif. Nous omettons la démonstration de sa
pleine fidélité.
\end{proof}

\begin{lemma}
\label{lemma-equivalence-properties}
Soient $X, x_i, U_i \to X, u_i$ comme dans
(\ref{equation-equivalence}). Si $f : Y \to X$ correspond aux
$g_i : Y_i \to U_i$ par $F$, alors $f$ est quasi-compact, quasi-séparé,
séparé, localement de présentation finie, de présentation finie, localement
de type fini, de type fini, propre, entier ou fini si et seulement si tous les
$g_i$ le sont, pour $i = 1, \ldots, n$.
\end{lemma}

\begin{proof}
Cela résulte de Limites d'espaces, lemme
\ref{spaces-limits-lemma-excision-modifications-properties}.
\end{proof}

\begin{lemma}
\label{lemma-equivalence-fibre}
Soient $X, x_i, U_i \to X, u_i$ comme dans
(\ref{equation-equivalence}). Si $f : Y \to X$ correspond aux
$g_i : Y_i \to U_i$ par $F$, alors
$Y_{x_i} \cong (Y_i)_{u_i}$ comme espaces algébriques.
\end{lemma}

\begin{proof}
C'est clair, puisque $u_i \to x_i$ est un isomorphisme.
\end{proof}







\section{Domination par des transformations quadratiques}
\label{section-quadratic-spaces}

\noindent
Nous ne définissons l'éclatement d'un espace en un point que si $X$ est
raisonnable.

\begin{definition}
\label{definition-blowup-at-point}
Soit $S$ un schéma. Soit $X$ un espace algébrique raisonnable sur $S$.
Soit $x \in |X|$ un point fermé. D'après Espaces raisonnables, lemme
\ref{decent-spaces-lemma-decent-space-closed-point}, on peut représenter $x$
par une immersion fermée $i : \Spec(k) \to X$.
Par {\it éclatement $X' \to X$ de $X$ en $x$}, on entend l'éclatement de $X$
le long du sous-espace fermé $Z = i(\Spec(k)) \subset X$.
\end{definition}

\noindent
À ce degré de généralité, l'éclatement de $X$ en $x$ n'est pas nécessairement
propre. Toutefois, si $X$ est localement noethérien, il résulte de
Diviseurs sur les espaces, lemme
\ref{spaces-divisors-lemma-blowing-up-projective}, que cet éclatement est
propre. Rappelons qu'un espace algébrique localement noethérien est noethérien
si et seulement s'il est quasi-compact et quasi-séparé. De plus, pour un
espace algébrique localement noethérien, être quasi-séparé équivaut à être
raisonnable (Espaces raisonnables, lemme
\ref{decent-spaces-lemma-locally-Noetherian-decent-quasi-separated}).

\begin{lemma}
\label{lemma-equivalence-sequence-blowups}
Soient $X, x_i, U_i \to X, u_i$ comme dans
(\ref{equation-equivalence}), et supposons que $f : Y \to X$ corresponde aux
$g_i : Y_i \to U_i$ par $F$. Il existe alors une factorisation
$$
Y = Z_m \to Z_{m - 1} \to \ldots \to Z_1 \to Z_0 = X
$$
de $f$ dans laquelle $Z_{j + 1} \to Z_j$ est l'éclatement de $Z_j$ en un
point fermé $z_j$ situé au-dessus de $\{x_1, \ldots, x_n\}$ si et seulement
si, pour chaque $i$, il existe une factorisation
$$
Y_i = Z_{i, m_i} \to Z_{i, m_i - 1} \to \ldots \to Z_{i, 1} \to Z_{i, 0} = U_i
$$
de $g_i$ dans laquelle $Z_{i, j + 1} \to Z_{i, j}$ est l'éclatement de
$Z_{i, j}$ en un point fermé $z_{i, j}$ situé au-dessus de $u_i$.
\end{lemma}

\begin{proof}
Un éclatement est un morphisme représentable. Dans l'un et l'autre cas,
on voit donc par récurrence que $Z_j \to X$ ou $Z_{i, j} \to U_i$ est
représentable. Par conséquent, chaque $Z_j$ ou $Z_{i, j}$ est un espace
algébrique raisonnable d'après Espaces raisonnables, lemme
\ref{decent-spaces-lemma-representable-named-properties}.
Cela montre que les assertions ont un sens (puisque l'éclatement n'est défini
que pour les espaces raisonnables).
Pour démontrer l'équivalence, partons d'une suite d'éclatements
$Z_m \to Z_{m - 1} \to \ldots \to Z_1 \to Z_0 = X$.
Le premier morphisme $Z_1 \to X$ est obtenu en éclatant l'un des $x_i$,
disons $x_1$. En appliquant $F$ à $Z_1 \to X$, on obtient l'éclatement
$Z_{1, 1} \to U_1$ en $u_1$, précisément l'éclatement en $u_1$, tandis que
$Z_{i, 0} = U_i$ pour $i > 1$.
À l'étape suivante, soit on éclate l'un des $x_i$, avec $i \geq 2$, sur
$Z_1$, soit on choisit un point fermé $z_1$ de la fibre de $Z_1 \to X$
au-dessus de $x_1$. Dans le premier cas, la marche à suivre est claire ;
dans le second, on utilise l'isomorphisme
$(Z_1)_{x_1} \cong (Z_{1, 1})_{u_1}$
(lemme \ref{lemma-equivalence-fibre}) pour obtenir le point fermé
$z_{1, 1} \in Z_{1, 1}$ correspondant à $z_1$. On définit alors
$Z_{1, 2} \to Z_{1, 1}$ comme l'éclatement en $z_{1, 1}$. En poursuivant
ainsi, on construit les factorisations de chacun des $g_i$.

\medskip\noindent
Réciproquement, étant données des suites d'éclatements
$Z_{i, m_i} \to Z_{i, m_i - 1} \to \ldots \to Z_{i, 1} \to Z_{i, 0} = U_i$
on construit exactement de la même manière la suite d'éclatements de $X$.
\end{proof}

\begin{lemma}
\label{lemma-make-ideal-principal}
Soit $S$ un schéma. Soit $X$ un espace algébrique noethérien sur $S$.
Soit $T \subset |X|$ un ensemble fini de points fermés $x$ tels que
(1) $X$ soit régulier en $x$ et (2) l'anneau local de $X$ en $x$ soit de
dimension $2$. Soit $\mathcal{I} \subset \mathcal{O}_X$ un faisceau
quasi-cohérent d'idéaux tel que $\mathcal{O}_X/\mathcal{I}$ soit supporté
par $T$. Il existe alors une suite
$$
X_m \to X_{m - 1} \to \ldots \to X_1 \to X_0 = X
$$
où $X_{j + 1} \to X_j$ est l'éclatement de $X_j$ en un point fermé $x_j$
situé au-dessus d'un point de $T$, telle que
$\mathcal{I}\mathcal{O}_{X_n}$ soit un faisceau inversible d'idéaux.
\end{lemma}

\begin{proof}
Écrivons $T = \{x_1, \ldots, x_r\}$. Choisissons des voisinages étales
élémentaires $(U_i, u_i) \to (X, x_i)$ comme dans la section
\ref{section-strategy}. Pour chaque $i$, la restriction
$\mathcal{I}_i = \mathcal{I}|_{U_i} \subset \mathcal{O}_{U_i}$
est un faisceau quasi-cohérent d'idéaux supporté en $u_i$.
L'anneau local de $U_i$ en $u_i$ est régulier et de dimension $2$.
Nous pouvons donc appliquer Résolution des surfaces, lemme
\ref{resolve-lemma-make-ideal-principal}, et obtenir une suite
$$
X_{i, m_i} \to X_{i, m_i - 1} \to \ldots \to X_1 \to X_{i, 0} = U_i
$$
d'éclatements en des points fermés situés au-dessus de $u_i$ telle que
$\mathcal{I}_i \mathcal{O}_{X_{i, m_i}}$ soit inversible.
Le lemme \ref{lemma-equivalence-sequence-blowups} fournit une suite
d'éclatements
$$
X_m \to X_{m - 1} \to \ldots \to X_1 \to X_0 = X
$$
comme dans l'énoncé du lemme, dont le changement de base aux $U_i$ redonne
les suites données. Il en résulte que $\mathcal{I}\mathcal{O}_{X_n}$ est un
faisceau inversible d'idéaux. En effet, nous le savons au-dessus de
$X \setminus \{x_1, \ldots, x_n\}$, et cela est vrai par construction dans
un voisinage étale de la fibre de chacun des $x_i$.
\end{proof}

\begin{lemma}
\label{lemma-dominate-by-blowing-up-in-points}
Soit $S$ un schéma. Soit $X$ un espace algébrique noethérien sur $S$.
Soit $T \subset |X|$ un ensemble fini de points fermés $x$ tels que
(1) $X$ soit régulier en $x$ et (2) l'anneau local de $X$ en $x$ soit de
dimension $2$. Soit $f : Y \to X$ un morphisme propre d'espaces algébriques
qui est un isomorphisme au-dessus de $U = X \setminus T$. Il existe alors
une suite
$$
X_n \to X_{n - 1} \to \ldots \to X_1 \to X_0 = X
$$
où $X_{i + 1} \to X_i$ est l'éclatement de $X_i$ en un point fermé $x_i$
situé au-dessus d'un point de $T$, ainsi qu'une factorisation
$X_n \to Y \to X$ de la composée.
\end{lemma}

\begin{proof}
D'après Compléments sur les morphismes d'espaces, lemme
\ref{spaces-more-morphisms-lemma-dominate-modification-by-blowup},
il existe un éclatement $U$-admissible $X' \to X$ qui domine $Y \to X$.
On peut donc supposer qu'il existe un faisceau d'idéaux
$\mathcal{I} \subset \mathcal{O}_X$ tel que
$\mathcal{O}_X/\mathcal{I}$ soit supporté par $T$ et que $Y$ soit
l'éclatement de $X$ le long de $\mathcal{I}$.
D'après le lemme \ref{lemma-make-ideal-principal}, il existe une suite
$$
X_n \to X_{n - 1} \to \ldots \to X_1 \to X_0 = X
$$
où $X_{i + 1} \to X_i$ est l'éclatement de $X_i$ en un point fermé $x_i$
situé au-dessus d'un point de $T$, telle que
$\mathcal{I}\mathcal{O}_{X_n}$ soit un faisceau inversible d'idéaux.
La propriété universelle de l'éclatement
(Diviseurs sur les espaces, lemme
\ref{spaces-divisors-lemma-universal-property-blowing-up})
fournit la factorisation voulue.
\end{proof}





\section{Domination par des éclatements normalisés}
\label{section-normalized-blowups}

\noindent
Dans cette section, nous montrons qu'une modification d'une surface peut être
dominée par une suite d'éclatements normalisés en des points.

\begin{definition}
\label{definition-normalized-blowup}
Soit $S$ un schéma. Soit $X$ un espace algébrique raisonnable sur $S$
satisfaisant aux conditions équivalentes de Morphismes d'espaces, lemme
\ref{spaces-morphisms-lemma-prepare-normalization}.
Soit $x \in |X|$ un point fermé. L'{\it éclatement normalisé de $X$ en $x$}
est la composée $X'' \to X' \to X$, où $X' \to X$ est l'éclatement de $X$
en $x$ (définition \ref{definition-blowup-at-point}) et où
$X'' \to X'$ est la normalisation de $X'$.
\end{definition}

\noindent
La normalisation $X'' \to X'$ est bien définie comme espace algébrique,
car $X'$ satisfait aux conditions équivalentes de
Morphismes d'espaces, lemme
\ref{spaces-morphisms-lemma-prepare-normalization}, d'après
Diviseurs sur les espaces, lemme
\ref{spaces-divisors-lemma-blowup-finite-nr-irreducibles}.
Voir Morphismes d'espaces, définition
\ref{spaces-morphisms-definition-normalization}
pour la définition de la normalisation.

\medskip\noindent
En général, l'éclatement normalisé n'est pas nécessairement propre, même
lorsque $X$ est noethérien. Rappelons qu'un espace algébrique est de Nagata
s'il admet un recouvrement étale par des espaces affines qui sont les spectres
d'anneaux de Nagata (Propriétés des espaces, définition
\ref{spaces-properties-definition-type-property} et
remarque \ref{spaces-properties-remark-list-properties-local-etale-topology},
et Propriétés, définition \ref{properties-definition-nagata}).

\begin{lemma}
\label{lemma-Nagata-normalized-blowup}
Dans la définition \ref{definition-normalized-blowup}, si $X$ est de Nagata,
alors l'éclatement normalisé de $X$ en $x$ est un espace algébrique normal
de Nagata, propre sur $X$.
\end{lemma}

\begin{proof}
Le morphisme d'éclatement $X' \to X$ est propre
(puisque $X$ est localement noethérien, nous pouvons appliquer
Diviseurs sur les espaces, lemme
\ref{spaces-divisors-lemma-blowing-up-projective}).
Ainsi, $X'$ est de Nagata (Morphismes d'espaces, lemme
\ref{spaces-morphisms-lemma-finite-type-nagata}).
La normalisation $X'' \to X'$ est donc finie (Morphismes d'espaces, lemme
\ref{spaces-morphisms-lemma-nagata-normalization}), et l'on en conclut que
$X'' \to X$ est également propre (Morphismes d'espaces, lemmes
\ref{spaces-morphisms-lemma-finite-proper} et
\ref{spaces-morphisms-lemma-composition-proper}).
Il s'ensuit que l'éclatement normalisé est un espace algébrique normal
(Morphismes d'espaces, lemme
\ref{spaces-morphisms-lemma-normalization-normal})
et de Nagata.
\end{proof}

\noindent
Voici, pour les éclatements normalisés, l'analogue du lemme
\ref{lemma-equivalence-sequence-blowups}.

\begin{lemma}
\label{lemma-equivalence-sequence-normalized-blowups}
Soient $X, x_i, U_i \to X, u_i$ comme dans
(\ref{equation-equivalence}), et supposons que $f : Y \to X$ corresponde aux
$g_i : Y_i \to U_i$ par $F$. Supposons que $X$ satisfasse aux conditions
équivalentes de Morphismes d'espaces, lemme
\ref{spaces-morphisms-lemma-prepare-normalization}.
Il existe alors une factorisation
$$
Y = Z_m \to Z_{m - 1} \to \ldots \to Z_1 \to Z_0 = X
$$
de $f$ dans laquelle $Z_{j + 1} \to Z_j$ est l'éclatement normalisé de $Z_j$
en un point fermé $z_j$ situé au-dessus de $\{x_1, \ldots, x_n\}$ si et
seulement si, pour chaque $i$, il existe une factorisation
$$
Y_i = Z_{i, m_i} \to Z_{i, m_i - 1} \to \ldots \to Z_{i, 1} \to Z_{i, 0} = U_i
$$
de $g_i$ dans laquelle $Z_{i, j + 1} \to Z_{i, j}$ est l'éclatement
normalisé de $Z_{i, j}$ en un point fermé $z_{i, j}$ situé au-dessus de $u_i$.
\end{lemma}

\begin{proof}
Cela résulte exactement du même argument que dans la démonstration du lemme
\ref{lemma-equivalence-sequence-blowups}.
\end{proof}

\noindent
Un espace algébrique de Nagata est localement noethérien.

\begin{lemma}
\label{lemma-dominate-by-normalized-blowing-up}
Soit $S$ un schéma. Soit $X$ un espace algébrique noethérien de Nagata sur
$S$, avec $\dim(X) = 2$. Soit $f : Y \to X$ un morphisme propre birationnel.
Il existe alors un diagramme commutatif
$$
\xymatrix{
X_n \ar[r] \ar[d] &
X_{n - 1} \ar[r] &
\ldots \ar[r] &
X_1 \ar[r] &
X_0 \ar[d] \\
Y \ar[rrrr]  & & & & X
}
$$
où $X_0 \to X$ est la normalisation et où $X_{i + 1} \to X_i$ est
l'éclatement normalisé de $X_i$ en un point fermé.
\end{lemma}

\begin{proof}
Bien que l'on puisse démontrer ce lemme directement pour les espaces
algébriques, nous poursuivrons la méthode employée ci-dessus pour nous ramener
au cas des schémas.

\medskip\noindent
Nous utiliserons le fait que les espaces algébriques noethériens sont
quasi-séparés et que leurs points possèdent donc des corps résiduels bien
définis (par exemple d'après Espaces raisonnables, lemme
\ref{decent-spaces-lemma-decent-space-elementary-etale-neighbourhood}).
Nous utiliserons sans autre mention les résultats de Morphismes d'espaces,
sections
\ref{spaces-morphisms-section-nagata},
\ref{spaces-morphisms-section-dimension-formula} et
\ref{spaces-morphisms-section-normalization}.
Nous pouvons remplacer $Y$ par sa normalisation. Soit $X_0 \to X$ la
normalisation. Le morphisme $Y \to X$ se factorise par $X_0$.
Nous pouvons donc supposer $X$ et $Y$ normaux.

\medskip\noindent
Supposons $X$ et $Y$ normaux. Le morphisme $f : Y \to X$ est un isomorphisme
au-dessus d'un ouvert qui contient tout point de codimension $0$ ou $1$ de
$Y$, ainsi que tout point de $Y$ au-dessus duquel la fibre est finie ; voir
Espaces au-dessus de corps, lemme
\ref{spaces-over-fields-lemma-modification-normal-iso-over-codimension-1}.
On voit donc qu'il existe un ensemble fini de points fermés $T \subset |X|$
tel que $f$ soit un isomorphisme au-dessus de $X \setminus T$.
D'après Compléments sur les morphismes d'espaces, lemme
\ref{spaces-more-morphisms-lemma-dominate-modification-by-blowup}
il existe un éclatement $X \setminus T$-admissible $Y' \to X$ qui domine $Y$.
Après avoir remplacé $Y$ par la normalisation de $Y'$, on voit que l'on peut
supposer le morphisme $Y \to X$ représentable.

\medskip\noindent
Écrivons $T = \{x_1, \ldots, x_r\}$. Choisissons des voisinages étales
élémentaires $(U_i, u_i) \to (X, x_i)$ comme dans la section
\ref{section-strategy}. Pour chaque $i$, le morphisme
$Y_i = Y \times_X U_i \to U_i$ est un morphisme propre birationnel qui est un
isomorphisme au-dessus de $U_i \setminus \{u_i\}$. Nous pouvons donc appliquer
Résolution des surfaces, lemme
\ref{resolve-lemma-dominate-by-normalized-blowing-up}
et obtenir une suite
$$
X_{i, m_i} \to X_{i, m_i - 1} \to \ldots \to X_1 \to X_{i, 0} = U_i
$$
d'éclatements normalisés en des points fermés situés au-dessus de $u_i$ telle
que $X_{i, m_i}$ domine $Y_i$.
Le lemme \ref{lemma-equivalence-sequence-normalized-blowups} fournit une suite
d'éclatements normalisés
$$
X_m \to X_{m - 1} \to \ldots \to X_1 \to X_0 = X
$$
comme dans l'énoncé du lemme, dont le changement de base à nos $U_i$ redonne
les suites données. Il résulte alors de l'équivalence de catégories du lemme
\ref{lemma-equivalence} que $X_m$ domine $Y$.
\end{proof}















\section{Changement de base au complété}
\label{section-aux}

\noindent
Le lemme simple qui suit sera un outil utile dans la suite.

\begin{lemma}
\label{lemma-iso-completions}
Soit $(A, \mathfrak m, \kappa)$ un anneau local dont l'idéal maximal
$\mathfrak m$ est de type fini. Soit $X$ un espace algébrique raisonnable
sur $A$. Posons $Y = X \times_{\Spec(A)} \Spec(A^\wedge)$, où $A^\wedge$
est le complété $\mathfrak m$-adique de $A$.
Soit $q \in |Y|$ un point d'image $p \in |X|$, situé au-dessus du point fermé
de $\Spec(A)$. L'homomorphisme d'anneaux locaux henséliens
$\mathcal{O}_{X, p}^h \to \mathcal{O}_{Y, q}^h$ induit un isomorphisme entre
leurs complétés.
\end{lemma}

\begin{proof}
Cela résulte immédiatement du cas des schémas : choisissons un voisinage
étale élémentaire $(U, u) \to (X, p)$ comme dans Espaces raisonnables, lemme
\ref{decent-spaces-lemma-decent-space-elementary-etale-neighbourhood},
et posons $V = U \times_X Y = U \times_{\Spec(A)} \Spec(A^\wedge)$ ainsi que
$v = (u, q)$.
Le cas des schémas est Résolution des surfaces, lemme
\ref{resolve-lemma-iso-completions}.
\end{proof}

\begin{lemma}
\label{lemma-port-regularity-to-completion}
Soit $(A, \mathfrak m, \kappa)$ un anneau local noethérien.
Soit $X \to \Spec(A)$ un morphisme localement de type fini, où $X$ est un
espace algébrique raisonnable. Posons
$Y = X \times_{\Spec(A)} \Spec(A^\wedge)$. Soit $y \in |Y|$
d'image $x \in |X|$. Alors :
\begin{enumerate}
\item si $\mathcal{O}_{Y, y}^h$ est régulier, alors
$\mathcal{O}_{X, x}^h$ est régulier ;
\item si $y$ appartient à la fibre fermée, alors
$\mathcal{O}_{Y, y}^h$ est régulier
$\Leftrightarrow \mathcal{O}_{X, x}^h$ est régulier ;
\item si $X$ est propre sur $A$, alors $X$ est régulier si et seulement si
$Y$ est régulier.
\end{enumerate}
\end{lemma}

\begin{proof}
Par localisation étale, les deux premières assertions résultent immédiatement
de l'analogue de ce lemme pour les schémas ; voir Résolution des surfaces,
lemme \ref{resolve-lemma-port-regularity-to-completion}.
Pour (3), puisque $Y \to X$ est surjectif (car $A \to A^\wedge$ est
fidèlement plat), la régularité de $Y$ entraîne celle de $X$ d'après (1).
Réciproquement, si $X$ est régulier, alors les anneaux locaux henséliens de
$Y$ sont réguliers en tous les points de la fibre spéciale.
Soit $y \in |Y|$ un point quelconque.
Dans le cas propre, l'application $|Y| \to |\Spec(A^\wedge)|$ est fermée ;
on peut donc trouver une spécialisation $y \leadsto y_0$, où $y_0$ appartient
à la fibre fermée. Choisissons un voisinage étale élémentaire
$(V, v_0) \to (Y, y_0)$ comme dans Espaces raisonnables, lemme
\ref{decent-spaces-lemma-decent-space-elementary-etale-neighbourhood}.
Comme $Y$ est raisonnable, on peut relever $y \leadsto y_0$ en une
spécialisation $v \leadsto v_0$ dans $V$ (Espaces raisonnables, lemme
\ref{decent-spaces-lemma-decent-specialization}).
On en conclut que $\mathcal{O}_{V, v}$ est un localisé de
$\mathcal{O}_{V, v_0}$, donc qu'il est régulier, ce qui achève la démonstration.
\end{proof}

\begin{lemma}
\label{lemma-formally-unramified}
Soit $(A, \mathfrak m)$ un anneau local noethérien. Soit $X$ un espace
algébrique sur $A$. Supposons que :
\begin{enumerate}
\item $A$ est analytiquement non ramifié
(Algèbre, définition \ref{algebra-definition-analytically-unramified}) ;
\item $X$ est localement de type fini sur $A$ ;
\item $X \to \Spec(A)$ est étale en tout point de codimension $0$ de $X$.
\end{enumerate}
Alors la normalisation de $X$ est finie sur $X$.
\end{lemma}

\begin{proof}
Choisissons un schéma $U$ et un morphisme étale surjectif $U \to X$.
Alors $U \to \Spec(A)$ satisfait aux hypothèses, donc aux conclusions, de
Résolution des surfaces, lemme \ref{resolve-lemma-formally-unramified}.
\end{proof}











\section{Propriétés qui s'en déduisent}
\label{section-existence-gives}

\noindent
Dans cette section, nous montrons que, pour un espace algébrique noethérien
intègre, l'existence d'une altération régulière entraîne de nombreuses
conséquences. Les lecteurs que n'intéressent pas les ``mauvais'' espaces
algébriques noethériens peuvent omettre cette section.

\begin{lemma}
\label{lemma-regular-alteration-implies}
Soit $S$ un schéma. Soit $Y$ un espace algébrique noethérien intègre sur $S$.
Supposons qu'il existe une altération $f : X \to Y$, avec $X$ régulier.
Alors la normalisation $Y^\nu \to Y$ est finie et $Y$ possède un ouvert dense
qui est régulier.
\end{lemma}

\begin{proof}
Par localisation étale, il suffit de le démontrer lorsque
$Y = \Spec(A)$, où $A$ est un domaine noethérien.
Soit $B$ la clôture intégrale de $A$ dans son corps des fractions.
Posons $C = \Gamma(X, \mathcal{O}_X)$. D'après
Cohomologie des espaces, lemme
\ref{spaces-cohomology-lemma-proper-pushforward-coherent}
nous voyons que $C$ est un $A$-module fini.
Comme $X$ est normal (Propriétés des espaces, lemme
\ref{spaces-properties-lemma-regular-normal})
nous voyons que $C$ est un domaine normal (Espaces au-dessus de corps, lemme
\ref{spaces-over-fields-lemma-normal-integral-sections}).
Ainsi $B \subset C$, et l'on conclut que $B$ est fini sur $A$ puisque $A$ est
noethérien.

\medskip\noindent
Il existe un ouvert non vide $V \subset Y$ tel que $f^{-1}V \to V$ soit fini ;
voir Espaces au-dessus de corps, définition
\ref{spaces-over-fields-definition-alteration}.
Quitte à restreindre $V$, on peut supposer que $f^{-1}V \to V$ est plat
(Morphismes d'espaces, proposition
\ref{spaces-morphisms-proposition-generic-flatness-reduced}).
Ainsi, $f^{-1}V \to V$ est fidèlement plat. Dès lors, $V$ est régulier d'après
Algèbre, lemme \ref{algebra-lemma-descent-regular}.
\end{proof}

\begin{lemma}
\label{lemma-regular-alteration-implies-local}
Soit $(A, \mathfrak m, \kappa)$ un domaine local noethérien.
Supposons qu'il existe une altération $f : X \to \Spec(A)$ avec $X$ régulier.
Alors :
\begin{enumerate}
\item il existe un $f \in A$ non nul tel que $A_f$ soit régulier ;
\item la clôture intégrale $B$ de $A$ dans son corps des fractions est finie
sur $A$ ;
\item le complété $\mathfrak m$-adique de $B$ est un anneau normal, c'est-à-dire
que les complétés de $B$ en ses idéaux maximaux sont des domaines normaux ;
\item la fibre formelle générique de $A$ est régulière.
\end{enumerate}
\end{lemma}

\begin{proof}
Les assertions (1) et (2) résultent du lemme
\ref{lemma-regular-alteration-implies}.
Il nous faut reprendre une partie de la démonstration de ce lemme afin de
fixer les notations pour la preuve de (3). Posons
$C = \Gamma(X, \mathcal{O}_X)$. D'après Cohomologie des espaces, lemme
\ref{spaces-cohomology-lemma-proper-pushforward-coherent}
nous voyons que $C$ est un $A$-module fini.
Comme $X$ est normal (Propriétés des espaces, lemme
\ref{spaces-properties-lemma-regular-normal})
nous voyons que $C$ est un domaine normal (Espaces au-dessus de corps, lemme
\ref{spaces-over-fields-lemma-normal-integral-sections}).
Ainsi $B \subset C$, et l'on conclut que $B$ est fini sur $A$ puisque $A$ est
noethérien. D'après Résolution des surfaces, lemme
\ref{resolve-lemma-algebra-helper}, il suffit, pour prouver (3), de montrer que
le complété $\mathfrak m$-adique $C^\wedge$ est normal.

\medskip\noindent
D'après Algèbre, lemme \ref{algebra-lemma-completion-finite-extension},
le complété $C^\wedge$ est le produit des complétés de $C$ en les idéaux
premiers de $C$ qui sont au-dessus de $\mathfrak m$. Ceux-ci sont en nombre
fini et ce sont
les idéaux maximaux $\mathfrak m_1, \ldots, \mathfrak m_r$ de $C$.
(Le résultat correspondant pour $B$ explique la dernière précision de
l'énoncé.) En remplaçant $A$ par $C_{\mathfrak m_i}$ et $X$ par
$X_i = X \times_{\Spec(C)} \Spec(C_{\mathfrak m_i})$, on se ramène donc au
cas du paragraphe suivant. (Notons que
$\Gamma(X_i, \mathcal{O}) = C_{\mathfrak m_i}$ d'après Cohomologie des espaces,
lemme \ref{spaces-cohomology-lemma-flat-base-change-cohomology}.)

\medskip\noindent
Ici, $A$ est un domaine local noethérien normal et $f : X \to \Spec(A)$ est
une altération régulière telle que $\Gamma(X, \mathcal{O}_X) = A$.
Il nous faut montrer que le complété $A^\wedge$ de $A$ est un domaine normal.
D'après le lemme \ref{lemma-port-regularity-to-completion},
$Y = X \times_{\Spec(A)} \Spec(A^\wedge)$ est régulier.
Or $\Gamma(Y, \mathcal{O}_Y) = A^\wedge$ d'après Cohomologie des espaces,
lemme \ref{spaces-cohomology-lemma-flat-base-change-cohomology}.
On conclut comme précédemment que $A^\wedge$ est normal.
En effet, $Y$ est normal d'après Propriétés des espaces, lemme
\ref{spaces-properties-lemma-regular-normal}.
Il est connexe parce que $\Gamma(Y, \mathcal{O}_Y) = A^\wedge$ est local.
Ainsi, $Y$ est normal et intègre (pour un espace algébrique séparé, être
connexe et normal entraîne être intègre). Par conséquent,
$\Gamma(Y, \mathcal{O}_Y) = A^\wedge$ est un domaine normal d'après
Espaces au-dessus de corps, lemme
\ref{spaces-over-fields-lemma-normal-integral-sections}.
Ceci démontre (3).

\medskip\noindent
Démontrons (4). Notons $\eta \in \Spec(A)$ le point générique et indiquons
par un indice $\eta$ le changement de base à $\eta$.
Comme $f$ est une altération, le schéma $X_\eta$ est fini et fidèlement plat
sur $\eta$. Puisque $Y = X \times_{\Spec(A)} \Spec(A^\wedge)$ est régulier
d'après le lemme \ref{lemma-port-regularity-to-completion}, on voit que
$Y_\eta$ est régulier (comme limite d'ouverts de $Y$).
Le morphisme $Y_\eta \to \Spec(A^\wedge \otimes_A \kappa(\eta))$ est alors
fini et fidèlement plat sur la fibre formelle générique. On conclut par
Algèbre, lemme \ref{algebra-lemma-descent-regular}.
\end{proof}











\section{Résolution}
\label{section-resolution}

\noindent
Voici une définition.

\begin{definition}
\label{definition-resolution}
Soit $S$ un schéma. Soit $Y$ un espace algébrique noethérien intègre sur $S$.
Une {\it résolution des singularités} de $X$ est une modification
$f : X \to Y$ telle que $X$ soit régulier.
\end{definition}

\noindent
Dans le cas des surfaces, nous souhaitons parfois disposer d'un peu plus
d'information.

\begin{definition}
\label{definition-resolution-surface}
Soit $S$ un schéma. Soit $Y$ un espace algébrique noethérien intègre de
dimension $2$ sur $S$. On dit que $Y$ admet une
{\it résolution des singularités par éclatements normalisés}
s'il existe une suite
$$
Y_n \to X_{n - 1} \to \ldots \to Y_1 \to Y_0 \to Y
$$
où :
\begin{enumerate}
\item $Y_i$ est propre sur $Y$ pour $i = 0, \ldots, n$ ;
\item $Y_0 \to Y$ est la normalisation ;
\item $Y_i \to Y_{i - 1}$ est un éclatement normalisé pour
$i = 1, \ldots, n$ ;
\item $Y_n$ est régulier.
\end{enumerate}
\end{definition}

\noindent
Remarquons que la condition (1) implique que la normalisation $Y_0$ de $Y$
est finie sur $Y$ et que les normalisations intervenant dans les éclatements
normalisés sont elles aussi finies.
Nous arrivons enfin au théorème principal de ce chapitre.

\begin{theorem}
\label{theorem-resolve}
Soit $S$ un schéma. Soit $Y$ un espace algébrique noethérien intègre de
dimension deux sur $S$. Les conditions suivantes sont équivalentes :
\begin{enumerate}
\item il existe une altération $X \to Y$ avec $X$ régulier ;
\item il existe une résolution des singularités de $Y$ ;
\item $Y$ admet une résolution des singularités par éclatements normalisés ;
\item la normalisation $Y^\nu \to Y$ est finie, $Y^\nu$ possède un nombre
fini de points singuliers $y_1, \ldots, y_m \in |Y|$, et, pour chaque $i$,
le complété de l'anneau local hensélien $\mathcal{O}_{Y^\nu, y_i}^h$ est
normal.
\end{enumerate}
\end{theorem}

\begin{proof}
Les implications (3) $\Rightarrow$ (2) $\Rightarrow$ (1) sont immédiates.

\medskip\noindent
Soit $X \to Y$ une altération avec $X$ régulier. Alors $Y^\nu \to Y$ est
fini d'après le lemme \ref{lemma-regular-alteration-implies}.
Considérons la factorisation $f : X \to Y^\nu$ de
Morphismes d'espaces, lemme \ref{spaces-morphisms-lemma-normalization-normal}.
D'après Espaces au-dessus de corps, lemme
\ref{spaces-over-fields-lemma-finite-in-codim-1}, le morphisme $f$ est fini
au-dessus d'un ouvert $V \subset Y^\nu$ contenant tout point de codimension
$\leq 1$ de $Y^\nu$. Alors $f$ est plat au-dessus de $V$ d'après
Algèbre, lemme \ref{algebra-lemma-CM-over-regular-flat}, et parce qu'un anneau
local normal de dimension $\leq 2$ est de Cohen--Macaulay par le critère de
Serre (Algèbre, lemme \ref{algebra-lemma-criterion-normal}).
L'ouvert $V$ est donc régulier d'après Algèbre, lemme
\ref{algebra-lemma-descent-regular}.
Comme $Y^\nu$ est noethérien, on conclut que
$Y^\nu \setminus V = \{y_1, \ldots, y_m\}$ est fini.
Pour chaque $i$, soit $\mathcal{O}_{Y^\nu, y_i}^h$ l'anneau local hensélien.
Alors $X \times_Y \Spec(\mathcal{O}_{Y^\nu, y_i}^h)$ est une altération
régulière de $\Spec(\mathcal{O}_{Y^\nu, y_i}^h)$
(nous omettons quelques détails).
D'après le lemme \ref{lemma-regular-alteration-implies-local}, le complété de
$\mathcal{O}_{Y^\nu, y_i}^h$ est normal. On obtient ainsi
(1) $\Rightarrow$ (4).

\medskip\noindent
Supposons (4). Il nous faut démontrer (3). Nous pouvons immédiatement
remplacer $Y$ par sa normalisation. Soient
$y_1, \ldots, y_m \in |Y|$ les points singuliers. Choisissons des voisinages
étales élémentaires $(V_i, v_i) \to (Y, y_i)$ comme dans la section
\ref{section-strategy}. Pour chaque $i$, l'anneau local hensélien
$\mathcal{O}_{Y^\nu, y_i}^h$ est le hensélisé de
$\mathcal{O}_{V_i, v_i}$. Ces anneaux ont donc des complétés isomorphes.
Le résultat pour les schémas (Résolution des surfaces, théorème
\ref{resolve-theorem-resolve}) montre alors qu'il existe des suites finies
d'éclatements normalisés
$$
X_{i, n_i} \to X_{i, n_i - 1} \to \ldots \to V_i
$$
dont les centres sont uniquement des points situés au-dessus de $v_i$, telles
que $X_{i, n_i}$ soit régulier. D'après le lemme
\ref{lemma-equivalence-sequence-normalized-blowups}, il existe une suite
d'éclatements normalisés
$$
X_n \to X_{n - 1} \to \ldots \to X_1 \to Y
$$
et, bien entendu, $X_n$ est lui aussi régulier (il suffit de considérer les
anneaux locaux). Cela achève la démonstration.
\end{proof}















\section{Exemples}
\label{section-examples}

\noindent
Voici quelques exemples liés aux résultats précédents de ce chapitre.

\begin{example}
\label{example-factorial}
\begin{reference}
\cite[4(c)]{Samuel-UFD}
\end{reference}
Soit $k$ un corps. L'anneau $A = k[x, y, z]/(x^r + y^s + z^t)$ est factoriel
lorsque $r, s, t$ sont des entiers deux à deux premiers entre eux.
En effet, puisque $x^r + y^s + z^t$ est irréductible, $A$ est un domaine.
L'élément $z$ est premier, c'est-à-dire qu'il engendre un idéal premier de
$A$. D'autre part, si $t = 1 + ers$ pour un certain $e$, alors
$$
A[1/z] \cong k[x', y', 1/z]
$$
où $x' = x/z^{es}$, $y' = y/z^{er}$ et $z = (x')^r + (y')^s$.
Ainsi, $A[1/z]$ est un localisé d'un anneau de polynômes, et il est donc
factoriel. Un argument de Nagata montre alors que $A$ est factoriel.
Voir Algèbre, lemme \ref{algebra-lemma-invert-prime-elements}.
Un argument analogue s'applique lorsque $t$ n'est pas congru à $1$ modulo
$rs$.
\end{example}

\begin{example}
\label{example-completion-not-factorial}
\begin{reference}
Voir \cite{Brieskorn} et \cite{Lipman-rational} pour la non-annulation, en
général, des groupes de Picard locaux.
\end{reference}
L'anneau $A = \mathbf{C}[[x, y, z]]/(x^r + y^s + z^t)$ n'est pas factoriel
lorsque $1 < r < s < t$ sont des entiers deux à deux premiers entre eux qui
ne sont pas égaux à $2, 3, 5$. Considérons par exemple le cas particulier
$A = \mathbf{C}[[x, y, z]]/(x^2 + y^5 + z^7)$.
Considérons les homomorphismes
$$
\psi_\zeta : \mathbf{C}[[x, y, z]]/(x^2 + y^5 + z^7) \to \mathbf{C}[[t]]
$$
donnés par
$$
x \mapsto t^7,\quad
y \mapsto t^3,\quad
z \mapsto -\zeta t^2(1 + t)^{1/7}
$$
où $\zeta$ est une racine $7$-ième de l'unité. Le noyau
$\mathfrak p_\zeta$ de $\psi_\zeta$ est un idéal premier de hauteur un ; si
$A$ était factoriel, il serait donc principal, engendré, disons, par
$f_\zeta \in \mathbf{C}[[x, y, z]]$.
Notons que $V(x^3 - y^7) = \bigcup V(\mathfrak p_\zeta)$ et que
$A/(x^3 - y^7)$ est réduit hors du point fermé. Par conséquent, toujours sous
l'hypothèse que $A$ est factoriel, on obtiendrait
$$
\prod\nolimits_\zeta f_\zeta = u(x^3 - y^7) + a(x^2 + y^5 + z^7)
\quad\text{in}\quad
\mathbf{C}[[x, y, z]]
$$
pour une unité $u \in \mathbf{C}[[x, y, z]]$ et un élément
$a \in \mathbf{C}[[x, y, z]]$. Après multiplication par une constante, on
peut supposer que $u(0, 0, 0) = 1$. Notons que le membre de gauche s'annule à
l'ordre $7$. Ainsi $a = - x \bmod \mathfrak m^2$. Mais il apparaît alors dans
le membre de droite un terme $xy^5$ qui n'apparaît pas dans celui de gauche.
C'est une contradiction.
\end{example}

\begin{example}
\label{example-not-blow-up}
Il existe un anneau local noethérien excellent de dimension $2$ et une
modification $X \to S = \Spec(A)$ qui n'est pas un schéma.
Esquissons une construction. Soit $X$ une surface normale sur $\mathbf{C}$
ayant un unique point singulier $x \in X$. Supposons qu'il existe une
résolution $\pi : X' \to X$ telle que la fibre exceptionnelle
$C = \pi^{-1}(x)_{red}$ soit une courbe projective lisse. Supposons en outre
qu'il existe un point $c \in C$ tel que, si $\mathcal{O}_C(nc)$ appartient à
l'image de $\Pic(X') \to \Pic(C)$, alors $n = 0$.
Soit ensuite $X'' \to X'$ l'éclatement au point non singulier $c$.
Notons $C' \subset X''$ la transformée stricte de $C$ et $E \subset X''$ la
fibre exceptionnelle. D'après les résultats d'Artin
(\cite{ArtinII} ; on peut utiliser, par exemple, \cite{Mumford-topology} pour
voir que le fibré normal de $C'$ est négatif), on peut contracter la courbe
$C'$ dans $X''$ et obtenir un espace algébrique $X'''$. Voici le diagramme :
$$
\xymatrix{
& X'' \ar[ld] \ar[rd] \\
X' \ar[rd] &  & X''' \ar[ld] \\
& X
}
$$
Nous affirmons que $X'''$ n'est pas un schéma. Cela fournit notre exemple,
car $X'''$ est un schéma si et seulement si le changement de base de $X'''$
à $A = \mathcal{O}_{X, x}$ est un schéma (nous omettons les détails).
Si $X'''$ était un schéma, l'image de $C'$ dans $X'''$ posséderait un voisinage
affine. Le complémentaire de ce voisinage serait un diviseur de Cartier
effectif sur $X'''$ (car $X'''$ est non singulier à l'exception de $1$ point).
Ce diviseur de Cartier effectif correspondrait à un diviseur de Cartier
effectif sur $X''$ qui rencontre $E$ et évite $C'$. En prenant son image dans
$X'$, on obtiendrait un diviseur de Cartier effectif qui rencontre $C$, au sens
ensembliste, en $c$. C'est impossible, puisque, par hypothèse, aucun multiple
de $c$ n'est la restriction d'un diviseur de Cartier.

\medskip\noindent
Pour terminer, il nous faut trouver une telle surface singulière $X$. On peut
simplement prendre pour $X$ la surface affine définie par
$$
x^3 + y^3 + z^3 + x^4 + y^4 + z^4 = 0
$$
dans $\mathbf{A}^3_\mathbf{C} = \Spec(\mathbf{C}[x, y, z])$, le point
singulier étant $(0, 0, 0)$. Alors $(0, 0, 0)$ est l'unique point singulier.
En éclatant $X$ le long de l'idéal maximal correspondant à $(0, 0, 0)$, on
obtient trois cartes, chacune isomorphe à la surface affine lisse
$$
1 + s^3 + t^3 + x(1 + s^4 + t^4) = 0
$$
qui est non singulière et dont le diviseur exceptionnel $C$ est défini par
$x = 0$. Le lecteur reconnaîtra en $C$ une courbe elliptique. Enfin, la
surface $X$ est rationnelle, comme le montre la projection depuis $(0, 0, 0)$,
ou encore parce que l'équation de l'éclatement permet de résoudre en $x$.
Le groupe de Picard d'une surface rationnelle non singulière est dénombrable,
tandis que celui d'une courbe elliptique sur les nombres complexes ne l'est
pas. On peut donc choisir un point fermé $c$ comme indiqué.
\end{example}








\begin{multicols}{2}[\section{Autres chapitres}]
\noindent
Préliminaires
\begin{enumerate}
\item \hyperref[introduction-section-phantom]{Introduction}
\item \hyperref[conventions-section-phantom]{Conventions}
\item \hyperref[sets-section-phantom]{Théorie des ensembles}
\item \hyperref[categories-section-phantom]{Catégories}
\item \hyperref[topology-section-phantom]{Topologie}
\item \hyperref[sheaves-section-phantom]{Faisceaux sur les espaces}
\item \hyperref[sites-section-phantom]{Sites et faisceaux}
\item \hyperref[stacks-section-phantom]{Champs}
\item \hyperref[fields-section-phantom]{Corps}
\item \hyperref[algebra-section-phantom]{Algèbre commutative}
\item \hyperref[brauer-section-phantom]{Groupes de Brauer}
\item \hyperref[homology-section-phantom]{Algèbre homologique}
\item \hyperref[derived-section-phantom]{Catégories dérivées}
\item \hyperref[simplicial-section-phantom]{Méthodes simpliciales}
\item \hyperref[more-algebra-section-phantom]{Compléments d'algèbre}
\item \hyperref[smoothing-section-phantom]{Lissage des morphismes d'anneaux}
\item \hyperref[modules-section-phantom]{Faisceaux de modules}
\item \hyperref[sites-modules-section-phantom]{Modules sur les sites}
\item \hyperref[injectives-section-phantom]{Objets injectifs}
\item \hyperref[cohomology-section-phantom]{Cohomologie des faisceaux}
\item \hyperref[sites-cohomology-section-phantom]{Cohomologie sur les sites}
\item \hyperref[dga-section-phantom]{Algèbre différentielle graduée}
\item \hyperref[dpa-section-phantom]{Algèbre à puissances divisées}
\item \hyperref[sdga-section-phantom]{Faisceaux différentiels gradués}
\item \hyperref[hypercovering-section-phantom]{Hyperrecouvrements}
\end{enumerate}
Schémas
\begin{enumerate}
\setcounter{enumi}{25}
\item \hyperref[schemes-section-phantom]{Schémas}
\item \hyperref[constructions-section-phantom]{Constructions de schémas}
\item \hyperref[properties-section-phantom]{Propriétés des schémas}
\item \hyperref[morphisms-section-phantom]{Morphismes de schémas}
\item \hyperref[coherent-section-phantom]{Cohomologie des schémas}
\item \hyperref[divisors-section-phantom]{Diviseurs}
\item \hyperref[limits-section-phantom]{Limites de schémas}
\item \hyperref[varieties-section-phantom]{Variétés}
\item \hyperref[topologies-section-phantom]{Topologies sur les schémas}
\item \hyperref[descent-section-phantom]{Descente}
\item \hyperref[perfect-section-phantom]{Catégories dérivées de schémas}
\item \hyperref[more-morphisms-section-phantom]{Compléments sur les morphismes}
\item \hyperref[flat-section-phantom]{Compléments sur la platitude}
\item \hyperref[groupoids-section-phantom]{Schémas en groupoïdes}
\item \hyperref[more-groupoids-section-phantom]{Compléments sur les schémas en groupoïdes}
\item \hyperref[etale-section-phantom]{Morphismes étales de schémas}
\end{enumerate}
Thèmes de théorie des schémas
\begin{enumerate}
\setcounter{enumi}{41}
\item \hyperref[chow-section-phantom]{Homologie de Chow}
\item \hyperref[intersection-section-phantom]{Théorie de l'intersection}
\item \hyperref[pic-section-phantom]{Schémas de Picard des courbes}
\item \hyperref[weil-section-phantom]{Théories cohomologiques de Weil}
\item \hyperref[adequate-section-phantom]{Modules adéquats}
\item \hyperref[dualizing-section-phantom]{Complexes dualisants}
\item \hyperref[duality-section-phantom]{Dualité pour les schémas}
\item \hyperref[discriminant-section-phantom]{Discriminants et différentes}
\item \hyperref[derham-section-phantom]{Cohomologie de de Rham}
\item \hyperref[local-cohomology-section-phantom]{Cohomologie locale}
\item \hyperref[algebraization-section-phantom]{Géométrie algébrique et formelle}
\item \hyperref[curves-section-phantom]{Courbes algébriques}
\item \hyperref[resolve-section-phantom]{Résolution des surfaces}
\item \hyperref[models-section-phantom]{Réduction semi-stable}
\item \hyperref[functors-section-phantom]{Foncteurs et morphismes}
\item \hyperref[equiv-section-phantom]{Catégories dérivées de variétés}
\item \hyperref[pione-section-phantom]{Groupes fondamentaux de schémas}
\item \hyperref[etale-cohomology-section-phantom]{Cohomologie étale}
\item \hyperref[crystalline-section-phantom]{Cohomologie cristalline}
\item \hyperref[proetale-section-phantom]{Cohomologie pro-étale}
\item \hyperref[relative-cycles-section-phantom]{Cycles relatifs}
\item \hyperref[more-etale-section-phantom]{Compléments de cohomologie étale}
\item \hyperref[trace-section-phantom]{La formule des traces}
\end{enumerate}
Espaces algébriques
\begin{enumerate}
\setcounter{enumi}{64}
\item \hyperref[spaces-section-phantom]{Espaces algébriques}
\item \hyperref[spaces-properties-section-phantom]{Propriétés des espaces algébriques}
\item \hyperref[spaces-morphisms-section-phantom]{Morphismes d'espaces algébriques}
\item \hyperref[decent-spaces-section-phantom]{Espaces algébriques décents}
\item \hyperref[spaces-cohomology-section-phantom]{Cohomologie des espaces algébriques}
\item \hyperref[spaces-limits-section-phantom]{Limites d'espaces algébriques}
\item \hyperref[spaces-divisors-section-phantom]{Diviseurs sur les espaces algébriques}
\item \hyperref[spaces-over-fields-section-phantom]{Espaces algébriques sur des corps}
\item \hyperref[spaces-topologies-section-phantom]{Topologies sur les espaces algébriques}
\item \hyperref[spaces-descent-section-phantom]{Descente et espaces algébriques}
\item \hyperref[spaces-perfect-section-phantom]{Catégories dérivées d'espaces}
\item \hyperref[spaces-more-morphisms-section-phantom]{Compléments sur les morphismes d'espaces}
\item \hyperref[spaces-flat-section-phantom]{Platitude sur les espaces algébriques}
\item \hyperref[spaces-groupoids-section-phantom]{Groupoïdes dans les espaces algébriques}
\item \hyperref[spaces-more-groupoids-section-phantom]{Compléments sur les groupoïdes dans les espaces}
\item \hyperref[bootstrap-section-phantom]{Amorçage}
\item \hyperref[spaces-pushouts-section-phantom]{Sommes amalgamées d'espaces algébriques}
\end{enumerate}
Thèmes de géométrie
\begin{enumerate}
\setcounter{enumi}{81}
\item \hyperref[spaces-chow-section-phantom]{Groupes de Chow des espaces}
\item \hyperref[groupoids-quotients-section-phantom]{Quotients de groupoïdes}
\item \hyperref[spaces-more-cohomology-section-phantom]{Compléments sur la cohomologie des espaces}
\item \hyperref[spaces-simplicial-section-phantom]{Espaces simpliciaux}
\item \hyperref[spaces-duality-section-phantom]{Dualité pour les espaces}
\item \hyperref[formal-spaces-section-phantom]{Espaces algébriques formels}
\item \hyperref[restricted-section-phantom]{Algébrisation des espaces formels}
\item \hyperref[spaces-resolve-section-phantom]{Résolution des surfaces revisitée}
\end{enumerate}
Théorie des déformations
\begin{enumerate}
\setcounter{enumi}{89}
\item \hyperref[formal-defos-section-phantom]{Théorie formelle des déformations}
\item \hyperref[defos-section-phantom]{Théorie des déformations}
\item \hyperref[cotangent-section-phantom]{Le complexe cotangent}
\item \hyperref[examples-defos-section-phantom]{Problèmes de déformation}
\end{enumerate}
Champs algébriques
\begin{enumerate}
\setcounter{enumi}{93}
\item \hyperref[algebraic-section-phantom]{Champs algébriques}
\item \hyperref[examples-stacks-section-phantom]{Exemples de champs}
\item \hyperref[stacks-sheaves-section-phantom]{Faisceaux sur les champs algébriques}
\item \hyperref[criteria-section-phantom]{Critères de représentabilité}
\item \hyperref[artin-section-phantom]{Axiomes d'Artin}
\item \hyperref[quot-section-phantom]{Espaces de Quot et de Hilbert}
\item \hyperref[stacks-properties-section-phantom]{Propriétés des champs algébriques}
\item \hyperref[stacks-morphisms-section-phantom]{Morphismes de champs algébriques}
\item \hyperref[stacks-limits-section-phantom]{Limites de champs algébriques}
\item \hyperref[stacks-cohomology-section-phantom]{Cohomologie des champs algébriques}
\item \hyperref[stacks-perfect-section-phantom]{Catégories dérivées de champs}
\item \hyperref[stacks-introduction-section-phantom]{Introduction aux champs algébriques}
\item \hyperref[stacks-more-morphisms-section-phantom]{Compléments sur les morphismes de champs}
\item \hyperref[stacks-geometry-section-phantom]{La géométrie des champs}
\end{enumerate}
Thèmes de théorie des espaces de modules
\begin{enumerate}
\setcounter{enumi}{107}
\item \hyperref[moduli-section-phantom]{Champs de modules}
\item \hyperref[moduli-curves-section-phantom]{Espaces de modules de courbes}
\end{enumerate}
Divers
\begin{enumerate}
\setcounter{enumi}{109}
\item \hyperref[examples-section-phantom]{Exemples}
\item \hyperref[exercises-section-phantom]{Exercices}
\item \hyperref[guide-section-phantom]{Guide bibliographique}
\item \hyperref[desirables-section-phantom]{Desiderata}
\item \hyperref[coding-section-phantom]{Style de programmation}
\item \hyperref[obsolete-section-phantom]{Obsolète}
\item \hyperref[fdl-section-phantom]{Licence de documentation libre GNU}
\item \hyperref[index-section-phantom]{Index généré automatiquement}
\end{enumerate}
\end{multicols}


\bibliography{my}
\bibliographystyle{amsalpha}

\end{document}
